%%%%%%%%%%%%%%%%%%%%%%%%%%%%%%%%%%%%%%%
% Cover Letter - Freeform, Additive Engineering Internship/Co-Op
% Compile with: pdflatex (Overleaf compatible)
%%%%%%%%%%%%%%%%%%%%%%%%%%%%%%%%%%%%%%%

\documentclass[11pt, a4paper]{article}

\usepackage[utf8]{inputenc}
\usepackage{mathptmx}           % Times New Roman
\usepackage[top=1in, bottom=1in, left=1.1in, right=1.1in]{geometry}
\usepackage{hyperref}
\usepackage{parskip}
\usepackage{setspace}
\usepackage{enumitem}

\hypersetup{
    colorlinks=true,
    urlcolor=blue,
    pdftitle={Mahima D. Prajapati - Cover Letter - Freeform},
}

\pagestyle{empty}

\begin{document}

% -------------------------------------------------------
%  HEADER
% -------------------------------------------------------
\begin{center}
    {\LARGE \textbf{Mahima D. Prajapati}}\\[4pt]
    Blacksburg, Virginia, USA \quad|\quad
    \href{mailto:mahimap@vt.edu}{mahimap@vt.edu} \quad|\quad
    \href{https://linkedin.com/in/mahimaprajapati}{linkedin.com/in/mahimaprajapati}
\end{center}

\vspace{8pt}
\hrule
\vspace{16pt}

% -------------------------------------------------------
%  DATE
% -------------------------------------------------------
\today

\vspace{16pt}

% -------------------------------------------------------
%  SALUTATION
% -------------------------------------------------------
Dear Freeform Hiring Team,

% -------------------------------------------------------
%  BODY
% -------------------------------------------------------
Freeform's mission to build AI-native, software-defined metal manufacturing systems maps directly onto the research I conduct every day. My PhD work at Virginia Tech addresses a fundamental challenge in directed energy deposition (DED): predicting and mitigating the structural distortions such as warping and cracking that arise from localized thermal gradients during the printing process. I am writing to apply for the Additive Engineering Internship/Co-Op for Summer or Fall 2026.

To tackle this problem, I am developing a process-aware, multiscale thermo-mechanical simulation framework in Abaqus that couples thermal history predictions with structural response models, material deposition sequencing, heat source modeling, and microstructure evolution at the pre-print, in-process, and post-print stages. The goal is a simulation tool that reflects the spatially and temporally varying material behavior that actually governs part performance in DED components, not idealized uniform properties. My three core contributions to date:

\begin{itemize}[leftmargin=1.5em, itemsep=2pt]
    \item \textbf{Validated thermo-mechanical FEA framework:} Thermal and structural models calibrated against published benchmarks; currently extending to microstructure evolution to move beyond assumed material uniformity.
    \item \textbf{ML surrogate for digital twin development:} Developing a physics-informed machine learning model trained on simulation outputs and process parameters to enable real-time process monitoring and predictive control, capabilities that high-fidelity FEA alone cannot provide at production speed.
    \item \textbf{Peer-reviewed publication:} MS research on aeroelastic analysis of wings with distributed electric propulsors published at AIAA SciTech 2026 (\href{https://doi.org/10.2514/6.2026-1649}{DOI: 10.2514/6.2026-1649}); awarded \textbf{First Place} at Virginia Tech's Graduate \& Professional Student Symposium.
\end{itemize}

The ML surrogate and digital twin thread of my research aligns precisely with what Freeform is building: a system where software closes the loop between process physics and part quality in real time. I want to contribute that perspective from the simulation side to the challenges Freeform faces on the factory floor.

I thrive in fast-paced environments where projects move from first principles to measurable outcomes quickly, and I would welcome the opportunity to bring that to your team.

I look forward to hearing from you.

% -------------------------------------------------------
%  CLOSING
% -------------------------------------------------------
\vspace{16pt}
Kind regards,

\vspace{32pt}
\textbf{Mahima D. Prajapati}

\end{document}
